% Generated draft: requires mathematical and source review.
\documentclass[11pt]{article}
\usepackage[margin=25mm]{geometry}
\usepackage{fontspec}
\setmainfont{Latin Modern Roman}
\usepackage{amsmath,amssymb,mathtools,unicode-math}
\setmathfont{Latin Modern Math}
\usepackage{enumitem,needspace,xurl,hyperref}
\setlength{\parindent}{0pt}
\setlength{\parskip}{5pt}
\setlength{\emergencystretch}{3em}
\newcommand{\N}{\mathbb{N}}
\newcommand{\Z}{\mathbb{Z}}
\newcommand{\Q}{\mathbb{Q}}
\newcommand{\R}{\mathbb{R}}
\newcommand{\C}{\mathbb{C}}
\newcommand{\F}{\mathbb{F}}
\newcommand{\A}{\mathbb{A}}
\newcommand{\PP}{\mathbb{P}}
\newcommand{\E}{\mathbb{E}}
\newcommand{\eps}{\varepsilon}
\newcommand{\dd}{\,\mathrm{d}}
\begin{document}
\section*{UnsolvedMath ID: 20001959}
\textbf{Source problem code:} catalogue entry
\subsection{Definitions and mathematical statement}
G. Landweber's Question 5.7 asks for an analogue of abstract moment maps that models local contact geometry. The source does not define the desired axioms or the exact local-model statement. The paired aid works with a compact Lie group \(G\) preserving a cooriented contact distribution \(\xi\) on \(M\).

Let \(\xi_+^\circ\subset T^*M\) be its positive annihilator cone, with tautological one-form \(\lambda\) and symplectic form \(d\lambda\). The reported intrinsic contact moment map is
\[
 \Psi:\xi_+^\circ\to\mathfrak g^*,\qquad
 \langle\Psi(\beta_x),\eta\rangle=\beta_x(\eta_M(x)),\qquad
 \Psi(r\beta_x)=r\Psi(\beta_x),\quad r>0.
\]
Use the Hamiltonian sign convention \(\iota_{\eta}\omega=-d\langle\Psi,\eta\rangle\). A \(G\)-invariant positive contact form \(\alpha\) defining \(\xi\) is a section of this cone, and the representative on the base is
\(\mu_\alpha=\alpha^*\Psi\), equivalently
\(\langle\mu_\alpha,\eta\rangle=\alpha(\eta_M)\).
Changing the form to \(f\alpha\), for smooth positive invariant \(f\), changes this representative to \(f\mu_\alpha\).

The aid records the exact pointwise kernel condition
\[
 \ker(\mu_\alpha(x))=
 \{\eta\in\mathfrak g:\eta_M(x)\in\xi_x\}.
\]
In particular isotropy Lie algebras are annihilated, and \(\mu_\alpha(x)|_{\mathfrak h}=0\) on \(M^H\). These are stronger contact-specific restrictions than allowing arbitrary constant projected values on fixed sets. Conversely the aid reports descent from homogeneous symplectic cones, with a free proper positive-dilation action and compatible \(G\)-action, to the corresponding cooriented contact quotient. Its full local-model hypotheses are not supplied.

For a fixed distribution and an invariant positive defining form \(\alpha_0\), the reported necessary-and-sufficient realizability test for a smooth equivariant map \(\nu\) is
\[
 \nu\text{ is induced by another invariant positive form defining }\xi
 \quad\Longleftrightarrow\quad
 \nu=f\mu_0\text{ for a smooth positive invariant }f,
 \qquad\mu_0=\mu_{\alpha_0}.
\]
The fixed-distribution and coorientation qualifications are essential.

The aid gives an explicit test case on standard contact \(S^3\subset\mathbb C^2\). For the circle action
\[
 e^{it}(z_1,z_2)=(e^{it}z_1,e^{-it}z_2),
\]
choose the standard invariant form normalized by
\(\mu_0=|z_1|^2-|z_2|^2\), and let \(\nu=\mu_0^3\). The aid calls \(\nu\) a GGK exact abstract moment map, obtainable from the invariant one-form \(\mu_0^2\alpha_0\). It has the same zero torus and the same positive rays away from zero as \(\mu_0\), yet any conformal multiplier off that torus is forced to be \(\mu_0^2\), which tends to zero there. Hence no everywhere positive smooth multiplier exists, and \(\nu\) is not contact-realizable in this fixed positive contact-form class. Abstract exactness alone therefore does not supply the stated contact realizability. The construction is partial progress on the requested general theory, not a complete axiomatization or universal local-normal-form theorem.
\subsection{Short English statement}
The aid formulates contact moment data intrinsically on the positive annihilator symplectic cone. On a fixed cooriented contact distribution, changing its defining form multiplies the representative moment map by a positive invariant function. A cubed moment map on the standard three-sphere is an abstract exact example that fails this fixed-distribution realizability test.
\subsection{Sources}
\begin{itemize}
\item Local source record: catalogue entry, catalogue entry; canonical \texttt{aim-geometry-notes.json}, \texttt{notes[296]}, source item 5.7 attributed to G. Landweber. Title: ``Intrinsic contact moment maps and a realizability sieve.''
\item AIM workshop list \emph{Moment maps and surjectivity in various geometries}, supplied reference \url{https://aimath.org/WWN/momentmaps/momentmaps.pdf}. Original remarks and literature fields are empty. The supplied verification repairs only line-break hyphenation ``mo-ment'' to ``moment''; this draft does not access a PDF.
\item Supplied research attempt 1 invokes the GGK exact abstract-moment-map framework and gives the cone description, realizability criterion and weight-\((1,-1)\) sphere example. No full GGK bibliographic entry or authoritative separate citation is supplied.
\item Registry dated 14 August 2026: ``partially\_solved,'' ``PARTIAL-PROGRESS,'' statement status ``corrected\_verified''; attempt partial, rigor ``proved,'' novelty confidence low. The machine-generated, non-peer-reviewed aid requires independent verification before citation.
\end{itemize}
\end{document}
